\section{Introduction}
Security operations centers must turn many detector outputs into a smaller number of reviewable incidents. Aggregation can reduce an analyst's workload, but a plausible relationship is not necessarily a trustworthy relationship: a copied claim may look like an independent witness, an event timestamp may precede the time its evidence became available, and a shared host or account may join unrelated activity. Research on domain-independent aggregation, multi-stage reconstruction, and triage addresses different parts of this problem \citep{landauer2022,wang2024,grain2025,sadlek2025}. The safety of a reconstructed incident depends on the evidence that justifies each cross-fragment union.

We study a specific decision problem: at a fixed cutoff, partition the alerts that have become visible into incident candidates, while refusing a cross-fragment link when its attestation is unavailable, dependent, inconsistent, or contradicted. The resulting groups are ranked for analyst review. Correlation proposes relationships; verification controls which extra unions can occur; ranking cannot repair a missing group. Prior work often emphasizes alert reduction, attack prediction, graph reconstruction, or prioritization, and these objectives and their labels are not interchangeable \citep{malach2025,grain2025,chuah2025}.

Our questions are: (RQ1) how does the implemented provenance and point-in-time verifier constrain cross-fragment reconstruction under incomplete evidence; (RQ2) does an endpoint-based fallback improve recoverability and retrieval on a matched, untouched synthetic holdout without observed unsafe unions; and (RQ3) what evidence limits further ranking gains and a valid external test? The contributions are an implementable account of the verifier and constrained fallback, a paired synthetic holdout comparison with group-safety measures, and a documented external data-readiness audit. We make no claim of superiority over published systems on incompatible datasets. Sections 2--5 establish the design; Sections 6--8 specify and report its evaluation; the remaining sections discuss interpretation and external validity.
